% !PRINTED 21
\weberpage{21}
%<<CONT
is set, there results from (4) the expression for the most general
algebraic differential belonging to the curve $f$:
\weq{5}{d\Omega = \frac{\Psi\Sigma \pm c_1x_2dx_3}{c_1f_1+c_2f_2+c_3f_3},}
in which $\Psi$ is an integral or fractional homogeneous function of
order $(n-3)$. The arbitrary quantities $c_1$, $c_2$, $c_3$
occur only apparently in this expression. In reality
it is entirely independent of them.

We consider again the case $n=3$. Then the simplest
case is that in which $\Psi$ is a constant, and (5) then passes over
into the elliptic differential of the first kind:
\weq{6}{du = \frac{\Sigma \pm c_1x_2dx_3}{c_1f_1+c_2f_2+c_3f_3}.}

This we now propose to transform, by means of the covariants of the ternary
form of the third degree, into the \gesperrt{Weierstrass} form.

In Bd.\,II, § 107 we became acquainted with the fundamental covariants of the
cubic form. These were, besides $f$ itself:
\[
\Delta = \frac{1}{6^3}\begin{vmatrix} f_{11}, & f_{12}, & f_{13} \\ f_{21}, & f_{22}, & f_{23} \\ f_{31}, & f_{32}, & f_{33} \end{vmatrix}, \quad
J = \frac{1}{36}\begin{vmatrix} f_{11}, & f_{12}, & f_{13}, & \Delta_1 \\ f_{21}, & f_{22}, & f_{23}, & \Delta_2 \\ f_{31}, & f_{32}, & f_{33}, & \Delta_3 \\ \Delta_1, & \Delta_2, & \Delta_3, & 0 \end{vmatrix},
\]
\[
K = \frac{1}{9}\begin{vmatrix} f_1, & \Delta_1, & J_1 \\ f_2, & \Delta_2, & J_2 \\ f_3, & \Delta_3, & J_3 \end{vmatrix},
\]
where the indices 1, 2, 3 denote differentiation with respect to the variables
$x_1$, $x_2$, $x_3$.

It follows from the multiplication theorem of determinants, since
$f$ and $df=0$, that:
\[
\begin{vmatrix} c_1, & c_2, & c_3 \\ x_1, & x_2, & x_3 \\ dx_1, & dx_2, & dx_3 \end{vmatrix}
\begin{vmatrix} f_1, & f_2, & f_3 \\ \Delta_1, & \Delta_2, & \Delta_3 \\ J_1, & J_2, & J_3 \end{vmatrix}
=
\begin{vmatrix} \Sigma x_i\Delta_i, & d\Delta \\ \Sigma x_iJ_i, & dJ \end{vmatrix}
\Sigma c_if_i
\]
\[
= 3\Sigma c_if_i(\Delta dJ - 2Jd\Delta),
\]
and from this
\weq{7}{du = 3\frac{\Delta dJ-2Jd\Delta}{K}.}
